\section{Actual finite products of original module localizations}

\begin{lemma}[A common denominator for the whole finite product]
\label{lem:finite-product-module-localization}
Let $I$ be finite, including the empty case, and let $a_i:M_i\to N_i$
be actual localizations at the SAME multiplicative subset $T$ of a
commutative ring $A$. Then the actual coordinate map
\[
 \prod_i M_i\longrightarrow\prod_i N_i,
 \qquad (x_i)_i\longmapsto(a_i(x_i))_i
\]
is localization at that SAME $T$.
\lean{EnvelopeProperty.FiniteProductModuleLocalization.piMap,EnvelopeProperty.FiniteProductModuleLocalization.piMap_apply,EnvelopeProperty.FiniteProductModuleLocalization.piMap_isLocalizedModule}
\leanok
\end{lemma}
\begin{proof}
Each element of $T$ acts invertibly on every actual factor, hence
on the whole product. Given an original tuple $(y_i)$, choose actual
numerators $x_i$ and denominators $s_i$ for its finitely many entries.
The product $s=\prod_i s_i$ is one denominator for the ENTIRE tuple:
multiply $x_i$ by the product of the remaining $s_j$. If two original
tuples have the same actual image, choose an annihilating denominator
for each coordinate and again take their product. This proves the
three native localization conditions for the actual coordinate map.
The empty products are one, and the empty tuple is unique, so no
nonempty-index assumption is needed.

For the whole principal-open section theorem this supplies only the
finite-product step. The actual affine section maps, coefficient
restrictions, Cech differential square and global kernel comparison
still must be constructed and are not replaced by a certificate.
\end{proof}
